% Generated draft: requires mathematical and source review.
\documentclass[11pt]{article}
\usepackage[margin=25mm]{geometry}
\usepackage{fontspec}
\setmainfont{Latin Modern Roman}
\usepackage{amsmath,amssymb,mathtools,unicode-math}
\setmathfont{Latin Modern Math}
\usepackage{enumitem,needspace,xurl,hyperref}
\setlength{\parindent}{0pt}
\setlength{\parskip}{5pt}
\setlength{\emergencystretch}{3em}
\newcommand{\N}{\mathbb{N}}
\newcommand{\Z}{\mathbb{Z}}
\newcommand{\Q}{\mathbb{Q}}
\newcommand{\R}{\mathbb{R}}
\newcommand{\C}{\mathbb{C}}
\newcommand{\F}{\mathbb{F}}
\newcommand{\A}{\mathbb{A}}
\newcommand{\PP}{\mathbb{P}}
\newcommand{\E}{\mathbb{E}}
\newcommand{\eps}{\varepsilon}
\newcommand{\dd}{\,\mathrm{d}}
\begin{document}
\section*{UnsolvedMath ID: 701406}
\textbf{Source problem code:} catalogue entry
\subsection{Definitions and mathematical statement}
Let \(\Gamma\) be a finite undirected graph with vertex set \(\{x_1,\ldots,x_n\}\), and write \(E(\Gamma)\) for its edges between distinct vertices. For group elements define \([a,b]=a^{-1}b^{-1}ab\). A group \(G\) is metabelian if its commutator subgroup \(G'\) is abelian, equivalently if \(G''=[G',G']=1\).

Let \(F=F(x_1,\ldots,x_n)\) be the free group, and define
\[
S_\Gamma=F\big/\left\langle\!\left\langle F''\cup\{[x_i,x_j]:\{x_i,x_j\}\in E(\Gamma)\}\right\rangle\!\right\rangle_F,
\]
where \(\langle\!\langle A\rangle\!\rangle_F\) denotes the normal subgroup of \(F\) generated by \(A\). Thus \(S_\Gamma\) is the partially commutative metabelian group specified in the source, with the edge commutation relations imposed in the variety \(\mathfrak A^2\) of all metabelian groups. This interprets the source's presentation conventionally as a presentation in that variety; nonedges are not additional defining relators. Loops, if permitted in the graph convention, contribute only tautological commutation relations.

In the first-order language of groups \(\mathcal L_{\mathrm{grp}}=\{\cdot,{}^{-1},1\}\), let
\[
\operatorname{Th}_{\forall}(S_\Gamma)=\left\{\forall y_1\cdots\forall y_m\,\varphi(y_1,\ldots,y_m):\ m\geq 0,\ \varphi\text{ is quantifier-free},\ S_\Gamma\models\forall y_1\cdots\forall y_m\,\varphi\right\}.
\]
Here quantifier-free formulas are Boolean combinations of equalities between group terms, with no parameters from \(S_\Gamma\).

\textbf{Question.} For every finite undirected graph \(\Gamma\), is \(\operatorname{Th}_{\forall}(S_\Gamma)\) decidable? That is, for each fixed \(\Gamma\), does there exist an algorithm which, given any universal \(\mathcal L_{\mathrm{grp}}\)-sentence, terminates and determines whether it holds in \(S_\Gamma\)? The source does not explicitly specify whether the algorithm is required to be uniform in \(\Gamma\), or whether an expanded language naming generators or other coefficients is intended; the formulation above uses the ordinary language of groups and asks decidability for each fixed graph.
\subsection{Short English statement}
For each finite undirected graph, form the group freely generated by its vertices subject to metabelianity and commutation along edges. Can an algorithm decide every universal sentence about this group in the ordinary language of groups? The source does not specify uniformity as the graph varies.
\subsection{Sources}
\begin{itemize}
\item Supplied UnsolvedMath local registry record: permanent ID \(701406\), exact source problem code \texttt{catalogue entry}, titled ``Kourovka Notebook Problem 17.104.'' The registry labels the problem \texttt{open}; this is supplied source metadata, not an independently verified mathematical or literature assessment.
\item \emph{The Kourovka Notebook (Unsolved Problems in Group Theory)}, 17th issue (2010), Problem 17.104; proposed by V. N. Remeslennikov and E. I. Timoshenko.
\item The supplied background also identifies the 21st issue edition (2026), \texttt{arXiv:1401.0300v48}. The collection description names E. I. Khukhro and V. D. Mazurov as editors. These bibliographic details are supplied attribution and have not been independently checked.
\end{itemize}
\end{document}
